% BEMv15 appendix snippet
% Conditional normalizer law for Route-B certified-length budget.
% This snippet is alpha-blind: it contains no observed fine-structure input.

\subsection{Conditional normalizer law for the certified Route-B correction}

\paragraph{Status.}
The result in this subsection is a conditional normalization certificate, not a
first-principles derivation of the fine-structure constant. It identifies the
minimal separable monomial normalizer compatible with the numerical BEMv14
evidence and with two structural assumptions on the truncated R--T spectral action.

Let
\[
S_M(K/0_1)
=
-\sum_{j=1}^M \log\lambda_j(K)
+
\sum_{j=1}^M \log\lambda_j(0_1)
\]
denote the finite pair-subtracted R--T spectral action after the BEM
discretization has produced \(M=M_{\max}\) retained modes. Let
\(L_{\rm cert}(K)\) denote the certified length from the ideal-knot database,
used only for the longitudinal scale.

\begin{assumption}[Mode extensivity]
For fixed geometry and fixed cutoff prescription, the unresolved finite
pair-action is extensive in the number of retained spectral degrees of freedom:
\[
\Delta F_{\rm pair}(K/0_1;M)
=
M\,\bar f(K/0_1)+o(M).
\]
Hence an intensive correction must contain a factor \(M^-1\).
\end{assumption}

\begin{assumption}[Length-scale covariance]
The finite R--T correction entering the certified longitudinal budget is a
second-order boundary density under global rescaling of the certified
longitudinal scale:
\[
L_{\rm cert}\mapsto \lambda L_{\rm cert},
\qquad
\bar f\mapsto \lambda^{-2}\bar f .
\]
Equivalently, the corresponding dimensionless intensive correction is obtained
by dividing the raw pair-action by \(M L_{\rm cert}^2\).
\end{assumption}

\begin{lemma}[Separable monomial normalizer]
Consider separable monomial normalizers of the form
\[
\mathcal N_{a,b}=M_{\max}^a L_{\rm cert}^b .
\]
Mode extensivity fixes \(a=1\). Length-scale covariance fixes \(b=2\). Thus
the minimal separable normalizer compatible with the two assumptions is
\[
\boxed{\mathcal N_{\rm RT}=M_{\max}L_{\rm cert}^2}.
\]
\end{lemma}

The certified Route-B budget then becomes
\[
\alpha^{-1}_{\rm pred,blind}
=
\frac12
\left[
N_{\rm soft}V_{\rm soft}L_{\rm cert}(3_1)
+
\frac{\Delta F_{\rm pair}(3_1/0_1)}
{M_{\max}L_{\rm cert}(3_1)^2}
\right].
\]
All quantities on the right are internal to the Route-B audit: no observed
fine-structure value is used.

\paragraph{Numerical certificate from BEMv14.}
Across the supplied BEMv14 multigrid stages the selected normalizer is
\[
\mathcal N_{\rm RT}=M_{\max}L_{\rm cert}^2 .
\]
The aggregated alpha-blind budget statistics are
\[
\langle \alpha^{-1}_{\rm pred,blind}\rangle
=
137.223563371624,
\qquad
\sigma
=
2.229777650702e-01,
\qquad
{\rm CV}
=
1.624923297367e-03.
\]
The largest normalized correction ratio in the supplied runs is
\[
\max
\frac{|\Delta F_{\rm pair}/(M_{\max}L_{\rm cert}^2)|}
{N_{\rm soft}V_{\rm soft}L_{\rm cert}}
=
9.304616834697e-03.
\]

\paragraph{Remaining open gate.}
The exponent \(b=2\) is not identifiable from single-knot data alone because
\(L_{\rm cert}(3_1)\) is fixed across the present grid. Its use is therefore
a scale-covariance assumption, not a purely empirical fit. A stronger closure
requires either a multi-knot certified-length test or an analytic heat-kernel
derivation of the length-squared density.
